\documentclass{article}
\usepackage{iclr2026_conference,times}

\usepackage{amsmath,amssymb}
\usepackage{booktabs}
\usepackage{graphicx}
\usepackage{microtype}
\usepackage{url}
\usepackage[table]{xcolor}
\usepackage{multirow}
\usepackage{hyperref}

\usepackage{xspace}
\sloppy

% ---- shared macros (keep every number in ONE place so text and tables cannot drift) ----
\newcommand{\nbank}{7{,}581}          % rules in the bank
\newcommand{\ntest}{1{,}170}          % clean test substrates
\newcommand{\nshared}{150}            % internal shared multi-method subset
\newcommand{\nglory}{37}              % external GLORYx substrates
\newcommand{\ceiling}{0.735}          % coverage_bank (micro)
\newcommand{\realised}{0.261}         % realised recall@15 (micro)
\newcommand{\selret}{0.368}           % selection retention
\newcommand{\grailmacro}{0.330}       % GRAIL macro recall@15
\newcommand{\interA}{+0.120}          % GRAIL<->BioTransformer interaction
\newcommand{\interAci}{[+0.073,\,+0.171]}
\newcommand{\interB}{+0.063}          % MetaTrans<->SyGMa interaction
\newcommand{\interBci}{[+0.022,\,+0.106]}

\title{The Bank Is Not the Bottleneck:\\
Decomposing Recall in\\
Metabolite-Structure Prediction}

% double-blind: authors stay anonymous until camera-ready
\author{Anonymous authors \\ Paper under double-blind review}

\newcommand{\fix}{\marginpar{FIX}}
\newcommand{\new}{\marginpar{NEW}}

% \iclrfinalcopy  % ENABLE ONLY for camera-ready -- it reveals authors and prints
%                 % "Published as a conference paper"; submission must stay anonymous.

\begin{document}
\maketitle

\begin{abstract}
Computational prediction of drug metabolites is assessed almost universally by top-$k$ recall, the
fraction of known metabolites a method recovers. A single number cannot say why a method scores
what it does, because a metabolite is missed for three unrelated reasons: no rule in the method's
knowledge base expresses the transformation, the applicable rule is not selected, or the correct
product is generated but ranked too low to report. Improving the rule base, the selector and the
ranker are different research programmes, and recall does not indicate which is warranted.

We decompose recall exactly into coverage, selection and ranking for methods built on an
enumerable rule base, and apply it to two of them. The comparison explains the leaderboard rather
than restating it. The method with the wider rule base loses: it reaches more chemistry and
converts a third of that reach, while its competitor reaches less and converts nearly all of it,
having no selection stage to lose anything at. The gap between them is the price of selecting at
all, which recall charges only to methods that select, and which no aggregate score reveals.

Every factor counts references matched against predictions, so the decomposition is defined only
relative to a stated rule for when two structures are the same molecule. That rule is never stated
in this field and has never been varied. We vary it over five conventions with predictions held
fixed, find the ordering of published methods is not invariant to the choice, and release the
protocol, the audited split and the frozen predictions that make the measurement repeatable.
\end{abstract}

\section{Introduction}
\label{sec:intro}

When a drug enters the body it is chemically transformed. Enzymes, principally in the liver,
oxidise, cleave and conjugate the administered compound into a set of new molecules called
metabolites. Which metabolites form is a central question in drug development, because a
metabolite may be inactive, may carry the therapeutic effect, or may be toxic in ways the parent
compound is not. Regulatory review attends closely to major human metabolites, and identifying them
experimentally is slow and expensive. Computational prediction of metabolite
structures therefore has direct practical value: given the structure of an administered compound,
predict the structures of the molecules the body will produce from it.

Two families of methods address this problem. Rule-based systems encode known biotransformations
as reaction templates, apply every applicable template to the input compound, and rank the
resulting products by an empirical score \citep{Ridder_2008, de_Bruyn_Kops_2020,
Djoumbou_Feunang_2019}. Machine-learning systems instead learn the transformation directly, either
as a sequence-to-sequence translation between molecular strings or as a classifier over reaction
types \citep{Litsa_2020, Zhu_2024}. Hybrid systems combine the two, using a learned model to
decide which of many templates to apply. Whatever the underlying approach, published methods are
compared on the same quantity: top-$k$ recall, the fraction of experimentally observed metabolites
that appear among a method's $k$ highest-ranked predictions.

That comparison hides what it is measuring. A rule-based or hybrid method can fail to recover a
metabolite for three unrelated reasons: the required transformation may have no corresponding
template in its rule base, the template may exist but not be selected, or the correct product may
be generated but ranked too low to be reported. Recall confounds these, so a low score identifies
no deficiency and suggests no remedy. Improving the rule base, improving template selection and
improving ranking are different research programmes, and the aggregate score does not indicate
which is warranted.

This work separates them. For a method whose rule base can be enumerated, recall factorises
exactly into coverage, selection and ranking, and the factorisation is a measurement rather than an
estimate: the coverage term is obtained by applying the entire rule base and asking what it
reaches. Applied to two such methods it explains a leaderboard instead of restating one. The method
with the wider rule base reaches more chemistry and converts a third of that reach; the one with
the narrower base converts nearly all of its own, because it applies every applicable rule and so
has no selection stage in which to lose anything. It wins by half a leaderboard position for that
reason alone.

Making the factorisation well-posed exposes a second problem, which is why the paper addresses it
first. Every factor counts references matched against predictions, so all three move with the rule
deciding when a predicted structure counts as the same molecule as a reported one, and chemistry offers
several defensible such rules: canonical SMILES strings, hashed chemical identifiers, identifiers
with stereochemistry removed, or molecular fingerprints. They disagree most about tautomers,
alternative structures of one compound differing in the position of a mobile hydrogen atom. Evaluations in this field adopt one each, usually
without stating why, and none has varied it. We vary it over five conventions with predictions held
fixed, and the ordering of published methods is not invariant to the choice. That a matching rule
can reorder a leaderboard has been shown before in other fields; that this field's rankings depend
on an unstated one has not.

We release the bundle the protocol needs, not a model: the leakage-audited split, the re-scoring
harness with all five criteria, the frozen predictions of all five methods under each, the
leakage-audit report, and GRAIL's checkpoints, GRAIL being the predictor we release so that the
decomposition has a second rule-grounded system to run on. Adding a sixth method needs its ranked
output on the split and nothing else.

The remainder of this article is organised as follows. Section~\ref{sec:related} reviews prior work
on metabolite prediction and on evaluation methodology in related fields. Section~\ref{sec:tame}
defines the matching criteria, the protocol and the audited split. Section~\ref{sec:decomp} states
the decomposition and the coverage ceiling. Section~\ref{sec:results} reports both, and
characterises the stage they identify as extreme multi-label classification.
Section~\ref{sec:limitations} discusses limitations and Section~\ref{sec:conclusion} concludes.

\section{Related Work}
\label{sec:related}

\paragraph{Metabolite-structure prediction.} Rule-based prediction of metabolite structures was
established by SyGMa \citep{Ridder_2008}, which encodes phase I (oxidation, reduction, hydrolysis) and phase II
(conjugation) biotransformations as reaction templates and scores the resulting products by the empirical frequency with which each
template yields an observed metabolite. GLORYx \citep{de_Bruyn_Kops_2020} couples a comparable
template library to a machine-learned model of sites of metabolism, the atoms at which an enzyme acts, and BioTransformer
\citep{Djoumbou_Feunang_2019} extends the template approach to microbial and environmental
transformations. Learned alternatives dispense with templates: MetaTrans \citep{Litsa_2020} treats
metabolite prediction as sequence-to-sequence translation between molecular strings, and
MetaPredictor \citep{Zhu_2024} combines a transformer with a filtering stage. These systems are
evaluated on overlapping but non-identical datasets and, as discussed below, under non-identical
notions of a correct prediction.

\paragraph{Comparative evaluations in this domain.} Several studies have placed these tools side by
side. \citet{Scholz_2023} benchmarked five predictors on 85 agrochemical parent compounds and
reported low first-generation precision together with substantial divergence between rule-based
and learned systems. \citet{Boyce_2022} compared six tools on 37 chemicals and found that the
system with the broadest coverage also produced the largest number of incorrect predictions.
\citet{Gao_2026} evaluated predictors against human radiolabelled absorption, distribution,
metabolism and excretion data and observed the same trade-off between coverage and precision. Each
of these studies fixes a single criterion for judging a prediction correct, and none reports how
sensitive its conclusions are to that criterion.

\paragraph{Structure standardisation.} The difficulty of deciding when two molecular structures
denote the same compound is documented independently of metabolite prediction.
\citet{Dhaked_2020} catalogue tautomeric transformations and show that the standard InChI
identifier normalises only a subset of them. \citet{Hahnke_2018} report that roughly 60\% of
PubChem records differ from their canonical InChI representation, predominantly through tautomer
choice, and \citet{Mansouri_2024} describe a standardisation pipeline developed for
quantitative structure-activity modelling. This literature treats standardisation as a
preprocessing step applied before analysis, rather than as a parameter of the evaluation itself.

\paragraph{Evaluation methodology in adjacent fields.} That methodological choices can alter the
apparent ordering of methods is recognised more broadly. \citet{Maziarz_2025} re-evaluated retrosynthesis models under a
standardised procedure and reversed a published ordering, partly because one model had been scored
under a relaxed stereochemistry criterion its baselines did not get: the closest precedent to this
work, and the only one that varies a structure-matching rule. \citet{agarwal2026metric} demonstrate for drug-response prediction that the choice
of scoring metric inverts which model performs best. \citet{liu2026massspecgymwild} audit published work built on a mass-spectrometry benchmark and
state the mechanism directly: divergent metric implementations, stereochemistry handling in
InChIKey matching among them, suffice to reorder rankings on identical predictions. Outside
chemistry the experiment has been run several times. re-scoring frozen text-to-SQL submissions under denotation
equivalence rather than exact match shows exact match misrepresenting the state of the art
\citep{Zhong_2020}, and re-scoring frozen question-answering outputs under four answer-equivalence
rules makes one system best under a rule it had lost under \citep{Kamalloo_2023}. Weighting evaluation examples by difficulty likewise reorders
natural-language leaderboards \citep{Mishra_2021, Rodriguez_2021}. Data splitting has received similar attention, and \citet{Joeres_2025}
propose a general procedure for constructing splits that minimise similarity between partitions.

\paragraph{Remaining gap.} Taken together, this literature establishes that structure
standardisation is unsettled, that evaluation choices can reorder method rankings in adjacent
domains, and that metabolite predictors have been compared without varying the criterion by which
a prediction is judged correct. That a matching rule can reorder a leaderboard is therefore not in doubt. What is unestablished is
whether it does so for metabolite predictors, and by how much, when the rule is the only thing
varying across a set of conventions rather than one relaxation of stereochemistry; and there is no
accepted way to attribute a predictor's score to a stage of its pipeline.

\section{TAME: an evaluation protocol with a declared matching criterion}
\label{sec:tame}

TAME is a \emph{protocol, an audited split and a re-scoring harness}, not a leaderboard service:
it exists so that the numbers in \S\ref{sec:results} are directly comparable and regenerable from
committed artifacts.

\paragraph{The matching criterion.} The field has never agreed what counts as a match. GLORYx uses a stereo-blind InChIKey skeleton, MetaTrans compares
Morgan fingerprints for identity, and rule-based systems that canonicalise their own output, such
as SyGMa, effectively compare canonical SMILES strings. Each is a different equivalence relation on molecules, and each is applied silently. TAME makes the choice explicit by exposing all of them as named quotients: predicted and
reference structures are mapped through one of
\[
q \in \{\textsc{canonical},\ \textsc{inchikey},\ \textsc{inchi-no-stereo},\ \textsc{tanimoto1},\
\textsc{inchikey-tautomer}\}
\]
and a prediction counts as a hit iff $q(\hat{y}) = q(y)$ for some reference $y$. Each map is
computed with RDKit 2022.09.5, the version that reproduces the columns reported here: \textsc{canonical} is canonical SMILES equality;
\textsc{inchikey} is the full InChIKey; \textsc{inchi-no-stereo} is its first block, the
connectivity skeleton, which is blind to stereochemistry, charge and isotopes; \textsc{tanimoto1}
is equality of Morgan fingerprints at radius 2 over 2048 bits, which is Tanimoto similarity of one;
and \textsc{inchikey-tautomer} standardises both sides before hashing, taking the fragment parent,
neutralising charges, canonicalising tautomers, and writing SMILES without stereochemistry.

Our recommended default is \textsc{inchikey-tautomer}. Standard InChI normalises only a
\emph{subset} of tautomers \citep{Dhaked_2020}, so a rule engine that emits a keto form (C=O) of
an enol reference (C=C-OH), two tautomers of the same compound, fails to match under plain
\textsc{inchikey} even though it produced the same chemical entity. Two properties of this default
must be stated because they bear on every number below. It is the most permissive of the five, and
it is stereo-blind as well as tautomer-blind, so a predicted D-alanine matches an L-alanine
reference. The stereo relaxation is not incidental: several methods here, GRAIL among them,
standardise their own output without stereochemistry and cannot emit a stereocentre at all, so a
stereo-sensitive criterion would score them on a representation they never produce. We therefore
report every result under all five criteria rather than defending one.

Unless stated otherwise all results use \textsc{inchikey-tautomer};
\S\ref{sec:rankflip} quantifies how much the ranking depends on that choice. A worked example in
which the same fixed predictions score a near-total miss or a full hit purely by changing $q$ is
given in Appendix~\ref{app:worked}.

Standardisation of this kind is established preprocessing (\S\ref{sec:related}); what is new
here is adopting it as the \emph{matching} rule and measuring how far it moves the leaderboard.

\paragraph{Split and leakage audit.} All splits are substrate-disjoint: no substrate in test or
validation appears in train, and no test substrate appears in validation. Clean triples are
constructed and then audited read-only by a separate script that canonicalises SMILES and verifies
zero substrate and zero positive-pair overlap across every split pair, emitting a machine-checkable
report. Leakage-aware splitting is itself not novel, since DataSAIL \citep{Joeres_2025} minimises cross-split similarity on
MoleculeNet-style tasks. The leak that matters here is
substrate--metabolite \emph{identity} overlap, which is domain-specific. As a secondary check, the deployed
configuration scored on both splits gives macro recall@15 $=0.344$ on validation and $0.334$ on
test. That evaluation covers 293 and 291 substrates respectively and is superseded on the full
split by the figures in Section~\ref{sec:diagnosis}. The gap of $+0.010$ lies in the expected
direction, validation being marginally easier because model selection touched it, but it is the
same magnitude as the seed-to-seed standard deviation reported in Section~\ref{sec:diagnosis}, so
it corroborates the zero-overlap audit rather than constituting independent evidence against
test-set overfitting.

\paragraph{Metrics.} We lead with recall@$k$ co-reported with the mean size of the output set, not with
precision alone. Under incomplete annotation precision is a \emph{pessimistic lower bound}: an
unannotated candidate is scored a false positive even when it is a real but unrecorded metabolite.
Reporting recall without output size is equally misleading in the other direction, since a method
can buy recall by enumerating (Table~\ref{tab:matchsens}, and Appendix~\ref{app:budget}). Model, preset and hyper-parameter
selection happens on validation, with two exceptions we name rather than imply. The deployed
filter's encoder variant was chosen by comparing both variants on test
(Appendix~\ref{app:arch}); that choice fixes which encoder is deployed and so reaches every GRAIL
number here, the headline included. The exploratory analyses of Appendix~\ref{app:negative} were
also scored on test, before this discipline was tightened, and support no claim anywhere. Results
we report on test --- the decomposition, the learned-versus-prior comparison, the breadth sweep
--- are measurements, not selections: no setting was chosen from them.

\section{The decomposition and the coverage ceiling}
\label{sec:decomp}

A rule-grounded predictor reaches a reference metabolite only by passing three gates in sequence:
the transformation must be
\emph{expressible} by some rule in GRAIL's bank of \nbank\ reaction rules, that rule must be
\emph{selected} into the substrate's candidate pool under a finite budget, and the resulting
product must be \emph{ranked} into the top $k$. Write, for a substrate set $S$ with reference sets
$Y_s$:
\[
\underbrace{\frac{|H|}{|U|}}_{\text{recall@}k}
\;=\;
\underbrace{\frac{|C_{\text{full}}|}{|U|}}_{\text{coverage}}
\;\cdot\;
\underbrace{\frac{|C_{\text{bud}}|}{|C_{\text{full}}|}}_{\text{selection}}
\;\cdot\;
\underbrace{\frac{|H|}{|C_{\text{bud}}|}}_{\text{ranking}},
\qquad
|U|=\textstyle\sum_{s\in S}|Y_s| ,
\]
where $C_{\text{full}}$ are references recovered by applying the \emph{entire} bank without a
budget, $C_{\text{bud}}$ those surviving the selector's budgeted pool, and $H$ those actually
returned in the top $k$. The identity is exact by construction (each set contains the next), so
the three factors partition the loss with no residual: it is an accounting device, not a model.
Its value is diagnostic: a single recall number cannot indicate
\emph{which} gate closed, while the three corresponding interventions, expanding the rule bank,
improving selection and improving ranking, address unrelated components of the pipeline. Two aggregations are in use, and every figure below is labelled with which one it is. The
\emph{micro} (ratio-of-sums) form pools reference metabolites across substrates, $|H|/|U|$ as
written above, so each substrate contributes in proportion to its number of references. The
\emph{macro} form averages per-substrate recall, weighting every substrate equally.
Table~\ref{tab:decomp} reports micro throughout, because only the micro form makes the three
factors multiply out exactly; the head-to-head comparisons in Section~\ref{sec:diagnosis} report
macro, which is the convention used by the comparator tools. On this split the two differ
($0.261$ micro against $0.330$ macro). Estimator details are in Appendix~\ref{app:formal}.

\paragraph{The ceiling as a primitive.} The first factor is measurable without any model: apply
all \nbank\ rules to every test substrate and ask what fraction of references appears anywhere in
the union. That is the bank's \textbf{coverage ceiling}, the best recall achievable if selection
and ranking were both perfect. On the clean test split it is \ceiling
(cluster-bootstrap 95\% CI $[0.709,0.762]$, resampling substrates). We argue this should be
reported alongside recall for any rule-grounded method, because it separates \emph{what the
knowledge base can express} from \emph{what the model does with it}, and because it bounds every
downstream claim: no ranking improvement can exceed it.

\begin{table}[t]
\caption{Coverage, selection and ranking factors for two rule-grounded methods on the clean test
split ($n=\ntest$), with cluster-bootstrap confidence intervals. Each row's factors multiply to
that method's realised recall@15.}
\label{tab:decomp}
\begin{center}
\begin{tabular}{lccccc}
\toprule
& coverage & selection & ranking & recall@15 & outputs \\
& $|C_{\text{full}}|/|U|$ & $|C_{\text{bud}}|/|C_{\text{full}}|$ & $|H|/|C_{\text{bud}}|$ & (micro) & per substrate \\
\midrule
GRAIL & \textbf{\ceiling} & \textbf{\selret} & 0.963 & \realised & 8.4 \\
      & \footnotesize$[0.709,0.762]$ & \footnotesize$[0.343,0.393]$ & \footnotesize$[0.946,0.978]$ & & \\
SyGMa & 0.542 & \textbf{1.000} & 0.942 & 0.511 & 81.7 \\
      & \footnotesize$[0.517,0.568]$ & \footnotesize\emph{by construction} & \footnotesize$[0.927,0.956]$ & & \\
\bottomrule
\end{tabular}
\end{center}
\end{table}

\paragraph{What the decomposition says.} Table~\ref{tab:decomp} localises the loss. The deployed
pipeline converts only $35.5\%$ of its own ceiling into realised recall, and the three factors
separate cleanly, with no pair of intervals overlapping. Selection is the binding constraint:
barely a third of what the bank can reach survives the selector's rule budget. Ranking is close to
lossless at $0.963$ $[0.946,0.978]$, so once a candidate is in the budgeted pool it almost always
reaches the output. The decomposition therefore points at selection, and specifically at how few
rules are applied, as the stage where intervention is warranted.

Applied to a second method it stops being self-diagnosis. SyGMa is rule-grounded with an
enumerable bank, so it decomposes on the same split in the same frame, and its profile is the
inverse of ours: its bank reaches less, $0.542$ against $\ceiling$, but it applies every applicable
rule and emits the result, so it converts almost all of that reach where we convert a third of a
wider one. The gap between us is therefore not a worse rule base but the cost of having a selection
stage at all, which recall charges only to methods that select. Two methods lose recall for
opposite reasons, and a single score says only that one of them scores lower. It also bounds what this architecture can achieve: even a perfect selector and a perfect
ranker would stop at \ceiling, and that ceiling is corpus-limited rather than a
consequence of tuning: 53\% of uncovered transformations are of reaction types absent from
the bank altogether (Appendix~\ref{app:props}).

\section{Results}
\label{sec:results}

\subsection{The criterion the decomposition rests on}
\label{sec:rankflip}

Because every factor of the decomposition counts matched references, the criterion has to be
settled before the decomposition means anything. Published methods have their \emph{frozen}
predictions re-scored by the TAME harness under all five criteria, so only the match rule changes
across columns, never the chemistry. Three of them
--- GRAIL, SyGMa \citep{Ridder_2008} and MetaPredictor \citep{Zhu_2024} --- have predictions for
the whole clean test split and are compared there, paired on identical substrates
(Table~\ref{tab:matchsens}). Two more, BioTransformer \citep{Djoumbou_Feunang_2019} and MetaTrans
\citep{Litsa_2020}, exist only on a 150-substrate shared subset; that five-method comparison is
reported in Appendix~\ref{app:xdomain}.

Methods are not equally sensitive to the convention. Over the full split GRAIL gains $0.016$
$[0.010,0.023]$ from canonical to tautomer matching against MetaPredictor's $0.040$
$[0.030,0.050]$, a differential of $+0.024$ $[0.014,0.034]$ surviving Holm--Bonferroni over the
three pairs. The same differential against BioTransformer on the 150-substrate subset is $+0.120$,
five times larger; we treat the full-split figure as the estimate, since a subset chosen for data
availability overstated the effect. Output size is a second, orthogonal confound: SyGMa emits $82$
candidates per substrate against GRAIL's $8.4$ (Appendix~\ref{app:budget}).



\begin{table}[t]
\caption{recall@15 by method and matching criterion on the full clean test split
($n=\ntest$), for the three methods with predictions on all of it. The rightmost column is each
method's gain from the strictest to the most tolerant criterion, with a paired-bootstrap interval.}
\label{tab:matchsens}
\begin{center}
\begin{tabular}{lccccc c}
\toprule
\textbf{method} & canon. & InChIKey & no-stereo & Tanimoto$=$1 & tautomer & gain (95\% CI) \\
 & & (strict) & (GLORYx) & (fingerprint) & (ours) & \\
\midrule
\textbf{GRAIL} & 0.323 & 0.331 & 0.333 & 0.324 & 0.340 & $+0.016$ $[0.010,0.023]$ \\
MetaPredictor  & 0.529 & 0.553 & 0.559 & 0.532 & 0.569 & $+0.040$ $[0.030,0.050]$ \\
SyGMa          & 0.529 & 0.554 & 0.559 & 0.531 & 0.567 & $+0.038$ $[0.029,0.048]$ \\
\bottomrule
\end{tabular}
\end{center}
\end{table}

\paragraph{Primary endpoint.} The endpoint is the \emph{differential} sensitivity, not any
individual reversal: the protocol $\times$ method interaction, how much more one method gains from
a more tolerant criterion than another, with a paired bootstrap over substrates ($10{,}000$
resamples; substrates are the independent unit). Every method gains under a more tolerant
criterion, so only the differential can reorder anything. The design was fixed before the data:
the five criteria were named, and the re-scoring of frozen predictions specified, twelve days
before the first result was computed. The estimator was not: we chose the differential and the
paired bootstrap when we analysed the results. We do correct for multiplicity. Under
Holm--Bonferroni at $\alpha=0.05$ every interaction this paper calls significant remains so, on
both sets; internally the two survive even when the family is taken conservatively as all ten
pairs the five-method table admits, rather than the two we computed. It is significant on two independent
pairs along two different protocol axes: BioTransformer gains \interA\ \interAci\ more than GRAIL
from canonical$\rightarrow$tautomer, and SyGMa gains \interB\ \interBci\ more than MetaTrans from
canonical$\rightarrow$InChIKey. The convention is therefore a \textbf{method-dependent
confounder}, not a neutral scoring choice.

\paragraph{Two reversals follow, neither individually significant.} Under canonical matching GRAIL
leads BioTransformer ($\Delta=+0.041$); under tautomer matching BioTransformer leads
($\Delta=-0.079$). MetaTrans leads SyGMa under canonical, no-stereo and Tanimoto$=$1 matching
($\Delta$ between $+0.009$ and $+0.013$) but trails under strict InChIKey ($\Delta=-0.053$), driven by a
\emph{non-monotone} response, since
MetaTrans is the only method whose recall \emph{falls} under
normalization, because it emits isomeric SMILES, which record the stereochemistry that a stereo-blind criterion discards and a strict one penalises. We note that every pairwise $\Delta$ here has a CI spanning zero at $n=\nshared$ (GRAIL$\leftrightarrow$
BioTransformer canonical $[-0.044,+0.122]$, tautomer $[-0.166,+0.007]$;
MetaTrans$\leftrightarrow$SyGMa canonical $[-0.064,+0.081]$, InChIKey $[-0.130,+0.023]$). The interaction underlying the reversals is therefore established, while the individual reversals are not.

\paragraph{External replication.} The results above live on a $\nshared$-substrate subset of our
own split, so we repeated the analysis on the field's de-facto external hold-out, the \nglory\
GLORYx drugs \citep{de_Bruyn_Kops_2020}, whose own tool is not scored here because the weights of
its site-of-metabolism component are not available for re-execution. Four methods are compared
across the criterion ladder (Table~\ref{tab:gloryx}). Sensitivity spans more than an order of
magnitude across them, and three of the four intervals exclude zero. Of the six method pairs, four show a significant interaction, all four surviving
Holm--Bonferroni over the six; the largest are GRAIL$\leftrightarrow$SyGMa ($-0.171$ $[-0.264,-0.086]$) and
MetaPredictor$\leftrightarrow$SyGMa ($-0.135$ $[-0.209,-0.072]$). The ordering consequence is
sharpest for the latter: under strict InChIKey SyGMa leads by a \emph{significant} $-0.130$
$[-0.196,-0.065]$; one rung of tolerance later, once stereochemistry is disregarded, that
advantage is gone and the point estimate reverses ($+0.005$, n.s.). A criterion chosen for convenience can therefore erase a
statistically significant advantage between two published tools. This pair is the sharpest of the six tested and survives
correction over all of them. This is shown on an external,
literature-standard set, with a different method quartet, at a sample size where the intervals
still exclude zero.

\begin{table}[t]
\caption{recall@15 on the external GLORYx set ($n=\nglory$) under three criteria of increasing
tolerance. The middle column separates the two relaxations that \textsc{inchikey-tautomer} applies
together; the interval is on the stereo step. The remaining tautomer step is $0.000$ for every
method except GRAIL, where it is $+0.023$ $[+0.000,+0.059]$.}
\label{tab:gloryx}
\begin{center}
\begin{tabular}{lcccl}
\toprule
\textbf{method} & InChIKey & no-stereo & tautomer & stereo step (95\% CI) \\
\midrule
SyGMa          & 0.492 & 0.498 & 0.498 & $+0.007$ $[+0.000,+0.020]$ \\
BioTransformer & 0.346 & 0.374 & 0.374 & $+0.028$ $[+0.002,+0.064]$ \\
MetaPredictor  & 0.362 & 0.504 & 0.504 & $+0.142$ $[+0.079,+0.215]$ \\
\textbf{GRAIL} & 0.155 & 0.310 & 0.333 & $+0.155$ $[+0.075,+0.244]$ \\
\bottomrule
\end{tabular}
\end{center}
\end{table}

\paragraph{Mechanism.} The spread is not arbitrary: a method's sensitivity tracks how far its
outputs diverge \emph{representationally} from the references it is scored against, and the
intermediate column identifies which divergence is responsible. On this set it is stereochemistry
alone: every gain that separates the methods is realised by the stereo step, and the tautomer step
that follows is zero for all but one method. GLORYx references carry stereocentres, and several of
these systems standardise their own output without stereochemistry and cannot emit one, so a
stereo-sensitive criterion scores them on a representation they never produce. Tautomerism is what
fails to match internally, where references are already stereo-stripped; stereochemistry is what
fails to match here. Which relaxation dominates is thus a property of the reference set, not of the
criterion. A single-model test in retrosynthesis,
a domain whose references are canonical and stereochemically explicit, finds a spread of about
$0.02$ across the same criteria, an order of magnitude below either effect
(Appendix~\ref{app:xdomain}).

\subsection{Comparison against published methods}
\label{sec:diagnosis}

GRAIL's several recall figures are not interchangeable; Appendix~\ref{app:xdomain} says which is which.
Before diagnosing GRAIL further we place it against the published methods. On the full clean test split ($n=\ntest$) the deployed GRAIL checkpoint reaches
\grailmacro\ macro recall@15; retraining the same configuration under three seeds gives
$0.344 \pm 0.010$ on that split. SyGMa reaches $0.572$, a paired difference of $-0.242$ from GRAIL
(95\% CI $[-0.271,\,-0.212]$, cluster bootstrap over $n=1{,}168$ common substrates).
MetaPredictor reaches $0.585$ on the shared subset ($n=\nshared$) and $0.568$ when re-scored on
the full test ($n=1{,}169$). GRAIL trails both. GRAIL does not achieve the highest recall at any output budget. A budget-matched frontier
sweeping $k \in [1,64]$ finds SyGMa ahead at \emph{every} $k$ with no crossover
(Appendix~\ref{app:budget}): the gap is not an artifact of output-size accounting, it is a genuine
selection weakness, and \S\ref{sec:decomp} localises it.

GRAIL differs from SyGMa on a second axis, the size of the returned set. On the same full test
split GRAIL returns $8.4$ candidates per substrate at precision $0.109$ (three-seed mean), where
SyGMa returns $81.0$ at $0.076$. This is an operating point that recall@$k$ alone does not
express.

\subsection{Selection as extreme multi-label classification}
\label{sec:xmc}

Since selection is the dominant shortfall, we ask what makes it hard, and the answer is structural
rather than metabolism-specific. Choosing which of \nbank\ rules apply to a substrate is an
\textbf{extreme multi-label classification} (XMC) problem with a near-empty positive set:
approximately one to three rules per
substrate carry a positive label, some ${\sim}0.03\%$ of the label space. Far more rules than that
are applicable to a given substrate; what is scarce is the annotation, not the applicability. In exactly
this regime XMC has long established that frequency/propensity baselines are hard to beat, which
is why propensity-scored losses were introduced \citep{Jain_2016}. Metabolite rule selection
\emph{inherits} this pathology, which the domain literature does not name.

This prediction is borne out. Each arm selects its own thirty rules and passes them through the
same downstream product enumeration and filter: a learned-only selector reaches $0.266$ while
selecting by empirical rule-firing frequency alone reaches $0.410$, a paired difference of
$-0.144$ (95\% CI $[-0.196,\,-0.095]$, $n=245$). The learned selector underperforms the frequency baseline. Under positive-unlabelled annotation with an approximately constant labelling
propensity, a learner with constant unlabelled weighting recovers a propensity-distorted score
whose dominant component is the marginal firing rate, so the prior is Bayes-competitive
\emph{by construction} at the selection stage; the learner can only win if substrate-conditional
prevalence variation exceeds estimation noise, and the observed data-scaling saturation indicates
it does not at this scale.

Two observations delimit the scope of this result. First, the effect
is \emph{selection-specific}: at the separate \emph{ranking} stage, on identical broad pools, the
learned filter beats the same frequency prior ($0.413$ versus $0.374$). Second, the pathology is
addressable rather than fundamental, since
a dense reformulation that predicts a coarse reaction type, an element-aware description of
which bonds change, rather than a specific rule exceeds the prior at retrieving the correct reaction
type within five guesses ($0.086$ against $0.030$; type recalls, not comparable with the
metabolite recalls elsewhere). That reformulation caps the fraction of references any selector could reach at about
47\%, so it diagnoses the problem rather than replacing the pipeline
(Appendix~\ref{app:props}).

\subsection{What constrains each factor}
\label{sec:evidence}

The diagnosis rests on seven measurements, not on the decomposition alone. Each one is designed to
move exactly one factor, so that a result which fails to move recall is as informative as one that
does. Table~\ref{tab:levers} collects them.

\begin{table}[t]
\caption{Seven diagnostic measurements, each attached to the factor of the decomposition it
constrains. Sample sizes differ between rows.}
\label{tab:levers}
\begin{center}
\small
\begin{tabular}{p{4.2cm}p{1.4cm}p{6.6cm}}
\toprule
\textbf{measurement} & \textbf{factor} & \textbf{result} \\
\midrule
learned selection against a frequency prior & selection &
$0.266$ against $0.410$; $\Delta = -0.144$ $[-0.196,-0.095]$ ($n{=}245$) \\
selection breadth, 30 to 300 rules & selection &
recall $0.352 \to 0.413$; pool coverage $0.482 \to 0.608$ ($n{=}245$) \\
training set, 2{,}418 to 4{,}787 substrates & selection &
recall $0.330 \to 0.334$, inside the $\pm0.010$ spread across seeds \\
\addlinespace
depth-2 rule application & coverage &
ceiling $+0.012$, at $8.5\times$ the candidate cost ($n{=}150$) \\
composition of the uncovered set & coverage &
$26.9\%$ of references outside the bank under \textsc{inchikey}, ${\sim}25\%$ under tautomer
matching; no mass-difference bin exceeds $6\%$ of the uncovered set ($n{=}500$) \\
\addlinespace
site-of-metabolism prior & ranking &
$+0.011$ at the best of five weightings, itself chosen on this split ($n{=}245$) \\
abstention on filter score & ranking &
precision never above $0.123$; recall $0.388 \to 0.041$ ($n{=}\ntest$) \\
\bottomrule
\end{tabular}
\end{center}
\end{table}

Four of the seven are interventions that failed, and they are the more useful half. Applying rules
twice, training on twice the data, adding a site-of-metabolism prior and abstaining on low filter
scores are the obvious remedies for a system in this position; none of them works here. Depth-2 application buys $0.012$ of ceiling for an $8.5$-fold increase in candidates, and
doubling the training set changes nothing, which is why we attribute the shortfall to the
structure of the selection problem rather than to under-training. Abstention is the sharpest of
the four: raising the threshold removes true and apparently false candidates in nearly equal
proportion, because an unannotated candidate is scored a false positive whether or not it is a real
metabolite, so precision stays flat while recall collapses. Precision here is bounded by the
annotation rather than by the threshold, which is why we report recall against output size rather
than precision. Of the remaining three, two
locate recoverable recall at the selection stage: the gap to the frequency prior and the gain from
widening selection. The third characterises what the bank cannot reach in one step.
Appendix~\ref{app:props} gives the measurements in full, with three propositions that account for
them.

\section{Limitations}
\label{sec:limitations}

\paragraph{Sample sizes differ.} The five-method table rests on $\nshared$
shared substrates, because two comparators could only be run on that subset; GRAIL and SyGMa are
scored on the full $\ntest$. The external replication is $n=\nglory$. At those sizes the
\emph{interactions} are significant but no individual pairwise reversal is, and we report them that
way throughout rather than promoting point estimates.

\paragraph{Two comparators cannot be re-run in this codebase.} MetaTrans's inference pipeline no longer
reproduces in our environment and BioTransformer's exact configuration is unscripted, so both are
represented by frozen prediction files rather than by re-executable code. Their numbers are
therefore reproducible from our artifacts but not independently regenerable from scratch, which is a weaker guarantee.

\paragraph{Precision is a lower bound, not an estimate.} Reference metabolite sets are
\emph{incomplete}: an unannotated but chemically real product is scored a false positive. Every
precision figure here should be read as pessimistic, which is why we lead with recall and output
size and never rank methods by precision.

\paragraph{The single-step bound.} Our ceiling is single-step-conditional. A non-vanishing fraction
of references are multi-generation products unreachable by any one rule application; depth-2
application lifts the ceiling by only ${\sim}0.012$ at approximately $8.5$ times the candidate
cost, so the $\ceiling$ ceiling bounds this single-step formulation rather than the task itself.

\paragraph{Evaluation hygiene of this paper itself.} Our headline numbers come from a single
evaluation of the deployed configuration, with model and hyper-parameter selection done on
validation. Some \emph{exploratory} analyses reported in Appendix~\ref{app:negative}, however, were
scored on the test split before we tightened this discipline. They inform no headline claim, and
their conclusions are negative. Subsequent exploratory work was moved to the validation split.

\section{Conclusion}
\label{sec:conclusion}

A recall figure for a rule-grounded predictor factorises, exactly and without further experiments,
into what its rule base can express, what its selector admits, and what its ranker returns. Applied
to two published systems the factorisation explains their ordering rather than repeating it: the
one with the wider rule base converts a third of what it reaches, the one with the narrower base
converts nearly all of its own because it selects nothing, and the gap between them is the price of
having a selection stage. Read as an aggregate, the same two numbers say only that one system
scores higher. We identify selection as the binding constraint for the predictor examined here, and
characterise that stage as extreme multi-label classification with very few positive labels per
instance, a regime in which frequency baselines are known to be hard to beat.

The factorisation is defined only relative to a rule for when two structures are the same molecule,
and that rule is a free parameter this field leaves unstated. It is not a neutral one: methods
differ systematically in how their outputs depart from reference structures, so varying it moves
the ordering, as it has been shown to move orderings in other fields. A reported recall figure is
therefore interpretable only alongside the criterion that produced it. Benchmark papers in this
area should report recall under all five criteria rather than fixing one silently, and, where the
method is rule-grounded, the coverage ceiling of its own rule base.

\subsubsection*{Reproducibility statement}

All numbers here are regenerable from committed artifacts, but for one reranker figure whose
per-substrate scores were not retained and which is flagged where it appears. The anonymized archive
accompanying this submission holds every method's frozen predictions, the leakage-audit report,
the harness with its five criteria, the split construction and audit, the analysis code, and the
GRAIL checkpoints.

\bibliographystyle{iclr2026_conference}
\bibliography{refs}

\appendix

\section{Architecture of GRAIL}
\label{app:arch}
GRAIL predicts xenobiotic metabolite structures in three stages. A learned generator scores every
rule in a bank of $\nbank$ SMIRKS templates against the substrate, estimating the probability
$P(r\mid s)$ that rule $r$ applies to substrate $s$. RDKit then applies each selected rule and
enumerates the products, recording which rule produced which candidate. A structural filter scores
each substrate--product pair, and the deployed system ranks candidates by the product of the
filter and generator scores, combining a judgement about the plausibility of the pair with the
generator's confidence in the rule that made it.

Both learned stages train on positive-unlabelled data. Annotated metabolites are the only
positives; rule-applicable products that carry no annotation are unlabelled, not negative,
because a missing annotation does not establish that a transformation does not occur. The filter
exposes its raw classifier outputs so that a non-negative positive-unlabelled surrogate is never
applied on top of a sigmoid when that estimator is selected, and the generator down-weights
applicable but unobserved rules instead of penalising them as hard negatives.

Both stages share a featurisation: molecular graphs with 16-dimensional nodes and 18-dimensional
edges, and a 1024-bit Morgan fingerprint per molecule. Unless stated otherwise, matching between
predicted and reference structures uses the tautomer-aware InChIKey convention throughout this
paper.

The three stages are separately inspectable: the selected rule, the enumerated product, and the
filter's judgement of the pair. What the architecture contributes is interpretable rule selection
paired with a positive-unlabelled filter, not recall above other metabolite predictors.

\subsection{The generator}

The generator scores the entire bank in one forward pass. A GATv2 encoder with three
message-passing layers, node dimensions $16\to192\to256\to128$, batch normalisation, ReLU and
dropout of 0.1 at each layer, and a mean-pooled readout, embeds the substrate graph into 128
dimensions. A parallel multilayer perceptron embeds its Morgan fingerprint, and a gated fusion
combines the two into the substrate representation. Each rule is embedded independently by a
rule-graph encoder of the same family, with hidden dimensions $128\to192\to128$, alongside a
six-dimensional branch of global rule features. The full bank is re-encoded on every forward pass,
which dominates the compute cost and is why inference time scales with bank size.

The logit for a rule sums four learned terms, each scaled by its own learned scalar. A
cross-attention term lets the rule embedding query the substrate's atom states through
single-head attention over node keys and values, producing a rule-specific local context that is
concatenated with the pooled embeddings and passed through a three-layer perceptron. A global term
is the cosine similarity between the substrate and rule embeddings, and a local term the cosine
similarity between the attention context and the rule. A fourth term reflects whether the rule's
SMARTS pattern matches. To this sum the model adds a per-rule bias and a frequency-prior term at
weight 0.4, the log-odds that the rule yields a true metabolite as estimated from training
statistics, then subtracts a penalty of 7.5 for rules whose SMARTS does not match, which
effectively masks inapplicable rules before the sigmoid. Where a rule fires at several sites, the
per-firing probabilities are combined by noisy-OR.

Training is multi-label rule classification with a margin-ranking auxiliary at weight 0.25 and
margin 0.45. Rules that match the substrate but carry no annotation are down-weighted by half
rather than treated as negatives, which is what makes the objective positive-unlabelled.
Optimisation uses Adam at a learning rate of $10^{-4}$ for at most 8 epochs with early stopping.

\subsection{The filter}

The deployed filter, which is the checkpoint behind every GRAIL number reported here, embeds the
substrate and the candidate product independently. Each passes through its own three-layer GATv2
encoder over the molecule's graph, with dimensions $16\to192\to256\to128$ and a mean-pooled
readout. The two 128-dimensional embeddings are concatenated with both Morgan fingerprints, and
the resulting 2304-dimensional vector is mapped through a three-layer perceptron
($\to128\to64\to1$, ReLU, dropout 0.1) to a single logit.

An alternative variant encodes one merged substrate--product graph, whose 18-dimensional nodes
are joined by cross-edges placed according to an element-aware maximum-common-substructure
correspondence between atoms, not by index order, so that chemically meaningful atom
mappings are preserved across the reaction. We measured the two rather than assuming which should
be the default, and report the comparison below.

The estimator is selectable: with the non-negative positive-unlabelled risk it consumes logits
and with binary cross-entropy it consumes probabilities, and the released configuration trains the
filter under the second. Optimisation uses Adam at a learning rate of $10^{-4}$ for at most 8 epochs
with early stopping at a patience of 7.

\subsection{Choosing between the two filters}

Both variants were trained on an identical 800-substrate subset, with the same negatives, seed and
optimiser, differing only in which encoder they used, and evaluated on the full clean test split
with the same generator and candidate pool.

Independent encoders win on accuracy. Recall@15 is 0.366 for the independent variant against 0.357
for the merged-graph variant, a paired difference of $-0.0090$ with a 95\% paired-bootstrap interval of
$[-0.017,-0.001]$ over $n=\ntest$ substrates; at recall@5 the margin widens, 0.282 against 0.259.

They also win on cost by roughly twenty-fold. Computing a maximum common substructure for every
pair makes full-data training take about 14 hours against roughly 40 minutes, and imposes the same
computation at every candidate at inference. Accuracy and cost point the same way. The merged-graph variant is implemented and available, and is not deployed because it loses on both counts.

\subsection{How the rule bank was built}

The rule-based system GRAIL is compared against, SyGMa, uses hand-authored reaction
rules. GRAIL's bank is mined automatically from annotated substrate--metabolite pairs, with each
candidate template validated against the pair that produced it.

Mining reads only the clean, molecule-disjoint training split. For each annotated pair it first
locates a ring-aware maximum common substructure, using RDKit's FMCS with ring atoms
matched only to ring atoms, rings completed, any bond order allowed to match, and atoms compared
by element; the match must cover at least 40\% of the smaller molecule's atoms. From that common
substructure it derives a reaction centre. Three kinds of atom are flagged: matched atoms whose element,
formal charge, hydrogen count, aromaticity or degree changed; atom pairs whose connecting bond
order changed; and leaving or entering atoms together with their matched neighbours. Several
substrate--product match combinations are tried, and the one giving the smallest centre is kept.
The centre is then
expanded by one bond to capture its immediate reactive environment, and the expanded fragments are
atom-mapped by their positional correspondence in the maximum common substructure into a candidate
SMIRKS template.

Two gates follow. A template survives only if applying it to its own source substrate regenerates
the annotated product after canonicalisation. The surviving, deduplicated templates then pass a
selectivity filter that rejects templates producing more than 200 distinct products on average
across a 50-substrate sample, and templates that are both very general and very prolific,
applicable to more than 90\% of sampled substrates while still averaging more than 50
products.

A second atom-mapping route is also implemented: it detects the reaction centre with RASCAL MCES,
falling back to FMCS, expands by one bond as above, and emits a mapped SMIRKS using RXNMapper, a
neural attention-based atom mapper. Both a rule-based and a neural mapping route are therefore
available.

Because mining reads the training split alone, no test metabolite can enter a mined rule, so the
coverage ceiling reported in the main text remains a held-out quantity.

\subsection{What the bank contains}

The deployed bank of $\nbank$ templates is the deduplicated union of four earlier curated banks,
of 473, 656, 500 and 1{,}051 rules, together with the newly mined templates.

Support across mined templates is heavily skewed. The positional miner alone extracts 5{,}856
unique self-tested templates from the training positives. Of these, 4{,}274 (73\%)
rest on a single training pair, and only 457 on five or more; the single best-supported template
covers 1{,}531 pairs. The bank is therefore dominated by highly specific, near-single-example
templates, not by broad general reactions.

Re-running the miner independently over the same training split re-derives templates that are all
already in the bank: 5{,}856 of 5{,}856, none new. Mining is deterministic and, on this corpus,
saturated at its source. The coverage ceiling of $\ceiling$ is bounded by the diversity of the
training reactions, not by incomplete mining, so extending coverage requires new reaction corpora
or more general templates rather than further mining of the same data.

The selectivity filter is the construction-time counterpart of the breadth--precision trade
measured in Appendix~\ref{app:props}. Rejecting over-general templates when the bank is built and
narrowing the number of rules applied at inference act on the same quantity, one before
deployment and one at the operating point.


\section{Formal framework}
\label{app:formal}
\paragraph{Overview.} The architecture described in the main text's architecture section, together with the coverage ceiling and recall decomposition reported later (in the coverage-ceiling and recall-decomposition sections), are unified by one compact generative model, which also gives the three propositions of the propositions section a principled home. The framework is deliberately thin: one likelihood, one decomposition, one lever$\to$factor map --- it is not a rewrite of the empirical results that follow.

\paragraph{A generative latent-reaction mixture.} For a substrate $s$, a metabolite $m$ arises by choosing a transformation rule $r$, choosing a firing site, and applying the rule:

\[
P(m \mid s) \;=\; \sum_{r} \sum_{\text{site} \,\in\, \text{sites}(r,s)} P(r \mid s) \cdot P(\text{site} \mid r, s) \cdot \mathbb{1}\big[\text{apply}(r, s, \text{site}) = m\big]
\]

The application term is deterministic given RDKit, so the inner indicator collapses to the enumerable product support; the latent variables are \emph{which rule fired} and \emph{at which site}, and the observation --- the annotated metabolite set --- is positive-unlabeled: a rule-applicable, non-annotated product is unlabeled, not a certified negative. The three deployed stages of the architecture section are one marginal-likelihood approximation of this model --- the model the pipeline approximates, not a claim that the trained weights are its MLE. The \textbf{generator} approximates $P(r \mid s)$ (rule selection); its persisted \texttt{rule\_prior\_logits} is the marginal $\pi(r) = P(r \text{ fires})$. \textbf{RDKit rule application} realises the deterministic support $\mathbb{1}[\text{apply}(r,s,\text{site})=m]$, enumerating candidate products. The \textbf{filter} approximates a discriminative correction $P(\text{true} \mid s, m)$ over the enumerated candidates. Deployment ranks by \texttt{filter\_score} $\times$ \texttt{generator\_score}. PU training approximates EM over the unobserved rule-firing indicator.

\paragraph{A recall decomposition.} Fix top-$k$. For substrate $s$, let $T_s$ be the true references and let the candidate sets nest $R_s(k) \subseteq P_{\text{bud},s} \subseteq P_{\text{full},s}$, where $P_{\text{full},s}$ is the full-bank depth-1 rule-applicable pool (the coverage ceiling reported in the coverage-ceiling section), $P_{\text{bud},s}$ the deployed budget-limited pool after generator selection, and $R_s(k)$ the top-$k$ ranked output. With $\text{hit}(A_s)$ the number of references matched by set $A_s$ under the tautomer-InChIKey quotient (monotone under inclusion) and pooled (micro) sums $U = \sum_s |T_s|$, $C_{\text{full}} = \sum_s \text{hit}(P_{\text{full},s})$, $C_{\text{bud}} = \sum_s \text{hit}(P_{\text{bud},s})$, $H = \sum_s \text{hit}(R_s(k))$:

\begin{align*}
\text{recall@}k &= \frac{H}{U} = \left(\frac{C_{\text{full}}}{U}\right) \cdot \left(\frac{C_{\text{bud}}}{C_{\text{full}}}\right) \cdot \left(\frac{H}{C_{\text{bud}}}\right) \\
&= \texttt{coverage\_bank} \cdot \texttt{selection\_retention} \cdot \texttt{ranking\_conversion}
\end{align*}

Because the sets nest, each factor lies in $[0,1]$ and the three multiply out exactly to the realised recall on any numbers --- it is an accounting identity, not a theorem; its only use is to localise where recall is lost across the pipeline, and its exact cancellation is never offered as evidence for anything beyond that bookkeeping. Qualitatively, \texttt{coverage\_bank} coincides with the rule-bank ceiling reported in the coverage-ceiling section.

\paragraph{Lever $\to$ factor map.} Every diagnostic in the propositions section and each proposition attaches to exactly one factor:

\begin{table}[h]
\centering
\caption{Lever $\to$ factor map: every diagnostic and proposition in the propositions section attaches to exactly one recall-decomposition factor.}
\label{tab:lever-factor-map}
\begin{tabular}{l p{9.5cm}}
\toprule
Factor & Levers \& propositions \\
\midrule
\texttt{coverage\_bank} & multi-step rule application (depth-2), the $\Delta$MW coverage gap, the external-validity composition covariate (see the external-validity section), Proposition 3 \\
\texttt{selection\_retention} & the learned-vs-prior rule-selection probe, data-scaling saturation, Proposition 2 \\
\texttt{ranking\_conversion} & the filter / listwise reranker, the oracle bound, Proposition 1 \\
\bottomrule
\end{tabular}
\end{table}


\section{Diagnosis of the recall shortfall}
\label{app:props}
\subsection{Diagnostic levers}

Each lever below attaches to exactly one factor of the decomposition introduced earlier (annotated in the Factor column); the table is followed by three refutable propositions that localise \emph{why}, not merely \emph{where}, recall is lost.

\begin{table}[h]
\centering
\footnotesize
\begin{tabular}{p{2.6cm}p{2.3cm}p{6.3cm}p{2.6cm}}
\toprule
Lever & Factor & Finding & Evidence \\
\midrule
\textbf{Learned vs.\ prior} \emph{(\texttt{top\_k}-limited rule-selection probe --- not deployed recall)} &
\texttt{selection\_\allowbreak retention} &
A SyGMa-style frequency prior significantly out-ranks the learned generator: gen-only recall@15 prior-only \textbf{0.410} vs learned-only \textbf{0.266} ($\Delta$ \textbf{$-$0.144}, 95\% CI \textbf{[$-$0.196, $-$0.095]}, paired bootstrap, n=245); with the filter, 0.405 vs 0.300 ($\Delta$ $-$0.105, CI [$-$0.152, $-$0.058]). Adding the prior to the learned scorer lifts \textbf{+0.130} gen / \textbf{+0.099} filter (both significant); the filter significantly helps only the \emph{weak learned} ordering (+0.034, CI [+0.011, +0.061]), not the strong prior (n.s.). &
\texttt{results/\allowbreak prior\_\allowbreak vs\_\allowbreak learned.json} (245 test subs) \\
\addlinespace
\textbf{Multi-step (depth-2)} &
\texttt{coverage\allowbreak\_bank} &
Breadth-capped depth-2 rule application lifts the ceiling by only \textbf{+0.012} (0.711$\to$0.723) at \textbf{8.5$\times$} candidate cost (194$\to$1653 candidates/substrate) --- not the dominant coverage lever. &
\texttt{results/\allowbreak benchmark\_\allowbreak report\_\allowbreak depth2.json} (150 subs, beam 10) \\
\addlinespace
\textbf{Coverage ($\Delta$MW gap)} &
\texttt{coverage\allowbreak\_bank} &
\textbf{26.9\%} of true metabolites are uncovered by the depth-1 bank (plain InChIKey; $\sim$25\% under tautomer). Misses are a diverse long tail --- the top missing class (hydroxylation/oxidation) is only \textbf{6\%} of uncovered. &
\texttt{results/\allowbreak benchmark\_\allowbreak report\_\allowbreak gap.json} (500 subs) \\
\addlinespace
\textbf{Data scaling} &
\texttt{selection\_\allowbreak retention} &
Recall saturates (2418$\to$4787 substrates $\approx$ flat) --- the plateau is not a data-quantity problem. &
\texttt{results/\allowbreak full\{2500,\allowbreak 5000\}\allowbreak \_single.log} \\
\addlinespace
\textbf{Regioselectivity (SoM)} &
\texttt{ranking\allowbreak\_conversion} &
A site-of-metabolism prior gives only a small lift --- regioselective ranking is hard within the bank. &
\texttt{results/\allowbreak train\_\allowbreak som.log} \\
\addlinespace
\textbf{Selection breadth (\texttt{top\_k} sweep)} &
\texttt{selection\_\allowbreak retention} &
Widening the deployed rule selector's \texttt{top\_k} 30$\to$300 lifts recall@15 \textbf{0.352$\to$0.413 (+0.061)} and pool coverage \textbf{0.482$\to$0.608}, saturating below the 0.735 ceiling (n=245 test subset; trend, not deployed-scale absolute). &
\texttt{results/\allowbreak selection\_\allowbreak ablation.json} (245 test subs) \\
\addlinespace
\textbf{Abstention (filter gate)} &
\texttt{ranking\allowbreak\_conversion} &
Gating output by a filter-score threshold $\tau$ traces a recall$\leftrightarrow$precision frontier, but precision stays \textbf{nearly flat ($\sim$0.10--0.12)} and no val-selected $\tau$ reaches 0.2: raising $\tau$ shrinks the output (9.96$\to$0.3) and collapses recall (\textbf{0.388$\to$0.041}) for essentially no precision gain --- precision is annotation-bounded, not threshold-bounded, justifying the rank-only default. &
\texttt{results/\allowbreak abstention\_\allowbreak frontier.json} (val-selected $\to$ 1170 test) \\
\bottomrule
\end{tabular}
\caption{Diagnostic levers and the recall-decomposition factor each one attaches to.}
\label{tab:levers}
\end{table}

\subsection{Selection breadth and abstention ablations}

\paragraph{Selection-breadth counterfactual (n=245 test subset).}
Sweeping the number of rules applied (\texttt{top\_k}) on the deployed pipeline (generator from \texttt{full5000\_priors} with the trained prior restored, deployed filter, ranked by the deployed \texttt{filter} $\times$ \texttt{generator} score) isolates rule-selection breadth as its own lever:

\begin{table}[h]
\centering
\begin{tabular}{lcccc}
\toprule
\texttt{top\_k} (rules applied) & pool coverage (oracle) & deployed recall@15 & mean pool & mean output \\
\midrule
30 (deployed) & 0.482 & 0.352 & 12.5 & 7.99 \\
100 & 0.547 & 0.388 & 35.3 & 9.47 \\
300 ($\approx$ applicable-rule limit) & 0.608 & 0.413 & 107.6 & 9.53 \\
\bottomrule
\end{tabular}
\caption{Selection-breadth counterfactual on the deployed pipeline (n=245 test subset).}
\label{tab:selection-breadth}
\end{table}

Widening \texttt{top\_k} 30$\to$300 monotonically lifts recall@15 \textbf{0.352$\to$0.413 (+0.061)} and pool coverage \textbf{0.482$\to$0.608}, at the cost of an \textbf{8.6$\times$} pool inflation (12.5$\to$107.6 mean candidates) --- recall bought with precision, the same trade SyGMa makes by overproducing ($\sim$74 outputs/substrate). It also saturates below the ceiling: pool coverage plateaus at \textbf{0.608 $<$ 0.735} even at $\sim$all applicable rules, and recall@15 0.413 still $\ll$ SyGMa's 0.572 --- so breadth is a real but partial lever. Ranking still loses roughly a third of in-pool hits even at \texttt{top\_k}=300 (recall / pool coverage = 0.413 / 0.608 = \textbf{0.68}, consistent with \texttt{ranking\allowbreak\_conversion} = 0.726, reported in the decomposition-results section).

\paragraph{Abstention is not a precision lever (val-selected $\tau$, full 1170-substrate test).}
Where breadth trades precision for recall, the reverse move --- deliberately abstaining --- barely moves precision. Gating the deployed pipeline's output by a filter-score threshold $\tau$ (candidates with \texttt{filter\_score} $<$ $\tau$ are dropped; survivors ranked by \texttt{filter} $\times$ \texttt{gen}, top-15) and sweeping $\tau$ traces a recall$\leftrightarrow$precision frontier whose precision axis is \textbf{nearly flat}: from rank-only $\tau=0$ (recall@15 \textbf{0.388}, precision \textbf{0.098}, 9.96 outputs/substrate) through $\tau=0.5$ (recall 0.146, precision 0.114, 1.55 outputs) to $\tau=0.9$ (recall 0.041, precision 0.040, 0.30 outputs), precision never exceeds \textbf{$\sim$0.123}, and no $\tau$ selected on validation reaches even precision 0.2. Raising $\tau$ removes true and apparently-false candidates in near-equal proportion --- because most ``false'' predictions are \emph{unannotated true} metabolites --- so gating shrinks the output and collapses recall for almost no precision gain; the best validation-selected F1 operating point ($\tau \approx 0.10$) lifts F1 only marginally (0.129$\to$0.131) by trimming the output tail. Precision here is therefore \textbf{annotation-bounded, not threshold-bounded}, which both substantiates the ``precision is a pessimistic lower bound under incomplete annotation'' caveat (noted in the evaluation-hygiene and decomposition-results sections) and justifies the rank-only deployment default (\texttt{results/\allowbreak abstention\_\allowbreak frontier.json}).

On the same broad pools, a second arm compares the deployed learned ranker (\texttt{filter} $\times$ \texttt{gen}) against the rule frequency prior alone (SyGMa-style):

\begin{table}[h]
\centering
\begin{tabular}{lcc}
\toprule
\texttt{top\_k} & filter $\times$ gen (learned) & frequency prior (SyGMa-style) \\
\midrule
30 & 0.352 & 0.359 \\
100 & 0.388 & 0.374 \\
300 & 0.413 & 0.374 \\
\bottomrule
\end{tabular}
\caption{Learned ranker vs.\ frequency prior on identical broad pools, by selection breadth.}
\label{tab:ranker-vs-prior}
\end{table}

At narrow breadth the two are $\approx$ equal (0.352 vs 0.359); as the pool broadens the \textbf{learned filter pulls ahead} (0.413 vs 0.374 at \texttt{top\_k}=300), while the frequency prior plateaus at 0.374 --- swapping to the prior signal at breadth would hurt, not help.

\emph{Synthesis:} GRAIL's deficit vs.\ SyGMa on this axis is a compound of (i) over-aggressive rule selection at the deployed \texttt{top\_k}=30 (recoverable $\sim$+0.06 by widening breadth) and (ii) a residual rule-set coverage gap (pool coverage saturates at 0.608 $<$ 0.735) --- it is \textbf{not} a ranking-model deficiency: the learned filter out-ranks the frequency prior once the candidate pool is broad. This is an \textbf{n=245 test-subset} ablation; the trend (breadth helps, ranking doesn't hurt) is the finding --- the absolute \texttt{top\_k}=30 baseline (0.352) is subset-inflated relative to the full-1170 deployed macro recall (0.330). (\texttt{results/\allowbreak selection\_\allowbreak ablation.json}, \texttt{results/\allowbreak selection\_\allowbreak ablation\_\allowbreak ranksignal.json})

\paragraph{Probe vs.\ deployed.}
The learned-vs-prior row is a controlled, \texttt{top\_k}-limited rule-selection \textbf{probe} (each mode picks its own top-30 rules, same downstream product loop and filter) --- it is \emph{not} the deployed recall. The deployed pipeline in fact applies its \textbf{top-30 selected rules} --- the selection-breadth ablation above confirms this: its \texttt{top\_k}=30 baseline (0.352 subset) reproduces deployed recall (0.330 macro / 0.344 3-seed, reported in the decomposition- and deployment-results sections), so the deployed pipeline \emph{is} a \texttt{top\_k}=30 rule selector, and widening that selector's breadth is itself a soft lever (previous paragraph), not evidence that rules are applied unselectively. Separately --- and narrower in scope --- the deployed pipeline is measurably \textbf{prior-independent} at that fixed \texttt{top\_k}=30 operating point: on an earlier 291-substrate evaluation, two checkpoints differing only in whether a trained prior buffer was present gave essentially identical recall (0.335 vs 0.334), so deployed recall (\textbf{0.330} macro / \textbf{0.261} micro, reported in the decomposition- and deployment-results sections) is coverage-limited, not prior-limited at the deployed operating point --- the prior \emph{term} barely moves recall there, distinct from the \texttt{top\_k} breadth axis above; the probe \emph{reveals} the learned selector's weakness; deployment \emph{masks} it by not depending on that selector's ranking. \emph{Honesty note:} an earlier draft reported the opposite --- ``learned beats prior'' --- an artifact of a checkpoint whose prior buffer was un-persisted; that reversal is withdrawn and corrected here, caught by adversarial verification. Neither multi-step application nor any single rule-family addition moves coverage much (misses are a diverse long tail), and data scaling is flat --- so the dominant, addressable loss remains \texttt{selection\_\allowbreak retention} (decomposition-results section: 0.489), the learned rule-selector performing worse than a trivial frequency prior in the probe. The three propositions below explain \emph{why}.

\subsection{Three propositions}

\paragraph{Proposition 1 --- Surrogate mismatch ($\to$ \texttt{ranking\allowbreak\_conversion}).}
A filter trained by a strictly proper scoring rule (BCE/PU) learns a globally calibrated posterior, Bayes-optimal for AUC and calibration; ranking each substrate's candidate pool by that posterior and taking top-$k$ is recall@$k$-optimal only when pools are homogeneous. GRAIL's pools vary in size (17--150) and positive rate (\texttt{n\_true} 1--18), so a pointwise-calibrated scorer can be recall@$k$-suboptimal even while a listwise, ranking-consistent surrogate dominates --- supported by a minimal 2-substrate counterexample (\texttt{grail\_\allowbreak metabolism/\allowbreak tests/\allowbreak test\_\allowbreak prop1\_\allowbreak counterexample.py}) in which the recall-superior reorder is verified \emph{not} globally calibrated, so a proper-scoring objective rejects it.

\emph{Confirmation --- a controlled objective swap on the full test set.} The cleanest evidence isolates the objective while holding everything else fixed: the same two auxiliary heads, the same frozen generator and filter, and the \textbf{same candidate pool} are used twice, differing only in how the heads were trained. Trained \textbf{pointwise} (independent MLE) and multiplied into the rank, they add \textbf{nothing} --- +0.0002 (95\% CI [$-$0.0073, +0.0077], n.s.). Fine-tuned instead with a \textbf{listwise ranking loss} against the frozen $\log(\text{filter})+\log(\text{gen})$ context, the identical heads gain \textbf{+0.0089 (95\% CI [+0.0032, +0.0153])} over their pointwise counterpart and \textbf{+0.0091 (95\% CI [+0.0027, +0.0160])} over the \texttt{filter} $\times$ \texttt{gen} baseline --- paired bootstrap over the full clean test, n=1170 (\texttt{results/\allowbreak joint\_\allowbreak rerank.json}). Since capacity, features, and pool are held constant, the margin is attributable to the \textbf{training objective}, which is exactly the surrogate mismatch this proposition predicts. A second, independent re-ranking arm reproduces the direction at the same n: \texttt{filter} $\times$ \texttt{gen} $\times$ \texttt{type} $\times$ \texttt{site} lifts recall@15 from 0.388 to 0.404, a paired \textbf{+0.0165 (95\% CI [+0.006, +0.027])}, while the factorized signal on its own is n.s. (+0.001, [$-$0.010, +0.013]) (\texttt{results/\allowbreak hybrid\_\allowbreak rerank\_\allowbreak full1170.json}).

\emph{Secondary, weaker evidence:} a higher-capacity listwise-InfoNCE reranker scored \textbf{0.433 $\to$ 0.500 @15 (+0.067)}, 74\% of the Stage-2 oracle 0.677, on a held-out Stage-2 run (\texttt{docs/\allowbreak benchmark/\allowbreak stage2\_\allowbreak ranker\_\allowbreak evidence.md}, Spike-3) --- reported as $\pm$0.015 across 3 seeds (pure reranker variance; the generator baseline is deterministic). We do \textbf{not} give this figure a paired interval: its per-substrate scores were not retained, so a paired bootstrap cannot be recomputed without retraining, and the proposition rests on the n=1170 intervals above rather than on this number. \emph{Guardrail:} this is a theorem about \textbf{objectives}, not a recall win --- the reranker's \textbf{0.500 still loses to SyGMa (0.558)} and is reported only as a separate Stage-2 artifact, never as a headline number.

\paragraph{Proposition 2 --- Learning loses to counting: rule selection is extreme multi-label classification ($\to$ \texttt{selection\_\allowbreak retention}).}
This is the paper's central selection finding, and it is not metabolism-specific folklore rediscovered: selecting which of the \textbf{7,581} rules apply to a substrate is an \emph{extreme multi-label classification} (XMC) problem with a near-empty positive set ($\approx$1--3 firing rules per substrate, $\sim$0.03\% of labels). In exactly this regime, XMC has long established that frequency/propensity baselines are hard to beat and that na\"{\i}ve losses are dominated by label sparsity --- propensity-scored losses were introduced precisely for this failure mode \citep{Jain_2016}. Metabolism rule selection \textbf{inherits this XMC pathology}, which the domain literature does not name; our contribution is to identify it here and to test it causally (below). Concretely: under PU annotation with an approximately constant labeling propensity (SCAR), a learner with constant unlabeled weighting recovers a propensity-distorted score whose dominant component is the marginal rule-firing rate $\pi(r)$ --- the frequency prior --- so the prior is Bayes-competitive \textbf{by construction} at the rule-selection stage (which rules clear \texttt{top\_k}), and the learned selector improves on it only if substrate-conditional prevalence variation exceeds estimation noise; the observed data-scaling saturation (Table~\ref{tab:levers}, row 4) indicates it does not at current scale.

\emph{Anchor:} learned-only \textbf{0.266} vs prior-only \textbf{0.410} (gen-only @15, $\Delta$ \textbf{$-$0.144}, 95\% CI \textbf{[$-$0.196, $-$0.095]}, paired bootstrap, n=245; \texttt{results/\allowbreak prior\_\allowbreak vs\_\allowbreak learned.json}). \emph{Ranking-stage nuance:} this prior advantage is selection-specific, not a general ``learned $<$ prior'' --- at the separate \textbf{ranking} stage (ordering the candidate pool once rules are already applied broadly), the learned filter is the better scorer: on the same broad pools, \texttt{filter} $\times$ \texttt{gen} beats the frequency prior \textbf{0.413 vs 0.374} at \texttt{top\_k}=300 (ranking-signal arm of the selection-breadth ablation above, n=245 subset; \texttt{results/\allowbreak selection\_\allowbreak ablation\_\allowbreak ranksignal.json}), consistent with Proposition 1's listwise-reranker gain. \emph{Falsifiable prediction:} reweighting the labeled loss by $1/\hat{e}(r)$ (a SAR correction) should shrink the prior's edge --- an \textbf{open test, not a promised fix}. \emph{Guardrail:} the propensity model $e(r) \propto \pi(r)$ is an \textbf{unmeasured modeling assumption}, flagged as such --- this is an explanatory model plus a refutable prediction, not proof that learning cannot win, and it is consistent with the deployed pipeline's prior-independence noted above.

\paragraph{Proposition 3 --- Paradigm limit ($\to$ \texttt{coverage\allowbreak\_bank}).}
Single-step rule-based recall $\leq$ single-step \texttt{coverage\allowbreak\_bank} $<1$, because a non-vanishing fraction of references are multi-generation (e.g.\ oxidation $\to$ conjugation) and unreachable by any single rule application --- an irreducible residual that no ranking improvement can recover.

\emph{Witnesses:} depth-2 rule application lifts the ceiling by only \textbf{+0.012} at \textbf{8.5$\times$} candidate cost (Table~\ref{tab:levers}, row 2; \texttt{results/\allowbreak benchmark\_\allowbreak report\_\allowbreak depth2.json}); on the external GLORYx set the uncapped single-step ceiling is \textbf{0.633} (reported in the external-validation section), whose references include single-step-unreachable multi-generation metabolites.

\emph{Composition of the gap:} typing each uncovered transformation by its radius-0 reaction type (element-aware changed-bond multiset) against the \textbf{4,417} types the bank already contains (\texttt{results/\allowbreak coverage\_\allowbreak gap\_\allowbreak types.json}; the pipeline reproduces the ceiling exactly, coverage 0.7355) shows the single-step gap is a mix skewed to novel chemistry --- of \textbf{687} uncovered test transformations, \textbf{41\% (280)} are of a reaction type the bank \emph{already} has (reachable in principle by a more general template, an in-bank ceiling of at most $(1910+280)/2597 = \textbf{0.843}$), while \textbf{53\% (366)} are \textbf{novel types absent from the bank}, addressable only by new reaction corpora (both together would bound the ceiling at 0.984). So template generalization is a real but \textbf{bounded} lever ($\approx$+0.11 ceiling headroom at best); the majority of the residual gap is genuinely out-of-bank chemistry, not a template-specificity artifact. \emph{Guardrail:} the bound is \textbf{single-step-conditional} --- it bounds the single-step paradigm, not the problem; multi-step and out-of-bank chemistry remain open coverage levers.

\subsection{Causal interventions}

\paragraph{Intervention: a factorized generator tests Propositions 2--3.}
This is a controlled intervention turning the observational decomposition above into a causal test: we rebuilt Stage-1 as a dense factorized generator, $P(\text{type}\mid s)\cdot P(\text{site}\mid \text{type},s)$, over a coarse $\sim$radius-0 type vocabulary.

\textbf{Selection is fixable:} the learned type head beats the frequency prior $\sim$3$\times$ at $k=5$ (recall@5 0.086 vs 0.030), robust across two vocabulary granularities (371 and 701 types), converging to the prior by $k=10$ (0.298 vs 0.309); site localisation holds (site hit@3 0.78, cf.\ Gate C 0.81) --- val set, \texttt{results/\allowbreak factorized\_\allowbreak val.json} --- so Proposition 2's PU degeneracy is a target-support artifact, not fundamental, confirming its falsifiable $1/\hat{e}$-style prediction.

\textbf{But coverage binds:} as a replacement generator the pipeline reaches only recall@15 0.256 on the full clean test --- a net regression against deployed GRAIL (budget-matched recall@15 0.365, precision 0.109) --- because type-selection caps oracle reachability at $\sim$47\%, unchanged when the vocabulary widened 371$\to$701 (\texttt{results/\allowbreak factorized\_\allowbreak eval.json}), confirming Proposition 3 causally.

\textbf{Ranking margin:} kept at broad coverage and used only to re-rank, \texttt{filter}$\times$\texttt{gen}$\times$\texttt{type}$\times$\texttt{site} lifts recall@15 from 0.388 to 0.404 over \texttt{filter}$\times$\texttt{gen} on the identical broad pool --- a paired +0.0165 (95\% CI [+0.006, +0.027], paired bootstrap over the full clean test, n=1170); the factorized signal alone is n.s.\ (+0.001, CI [$-$0.010, +0.013]) (\texttt{results/\allowbreak hybrid\_\allowbreak rerank\_\allowbreak full1170.json}).

\textbf{Joint vs.\ bolt-on ranking:} the bolt-on re-ranker multiplies \emph{independently} MLE-trained type/site probabilities into the rank; fine-tuning those two heads instead with a listwise ranking loss against the \textbf{fixed} generator and filter (frozen $\log(\text{filter})+\log(\text{gen})$ context, gradients only through the heads) sharpens the signal --- on a matched \texttt{top\_k}=100 pool the joint-trained re-ranker beats the bolt-on by a paired \textbf{+0.0089 (95\% CI [+0.003, +0.015], n=1170)} and the \texttt{filter}$\times$\texttt{gen} baseline by \textbf{+0.0091 [+0.003, +0.016]}, where on that pool the independently-trained bolt-on adds nothing (+0.0002, n.s.). So \emph{rank-aware} training, not the type/site heads per se, recovers this margin (\texttt{results/\allowbreak joint\_\allowbreak rerank.json}).

\emph{Synthesis:} selection is addressable, ranking yields a small significant margin, but coverage --- the source-limited rule bank --- remains binding; no Stage-1 redesign matches SyGMa without a broader bank.

\paragraph{Set-generation (GFlowNet): a released negative result.}
Beyond re-ranking, we also formulated Stage-1 as a \textbf{set-terminal GFlowNet} (\texttt{model/\allowbreak set\_\allowbreak gflownet.py}): the terminal object is a forest/set of metabolites per substrate, trained by Trajectory Balance against a PU set-coverage reward with an analytic \texttt{1/\#leaves} backward --- a genuinely novel formulation for this domain, released for reuse. On our benchmark it does \textbf{not} beat the simpler baselines: across 3 seeds (val, n=40) its single-set recall@15 is \textbf{0.292$\pm$0.032}, statistically indistinguishable from beam (0.311$\pm$0.039) and below the Stage-2a reranker (0.411$\pm$0.044). It does produce diverse sets (pairwise-Tanimoto 0.214$\pm$0.003, $\sim$39 unique scaffolds), but a clean \emph{diversity-advantage} claim is \textbf{not} supported here --- the union-coverage comparison against the baselines is dominated by an under-production artefact in the baseline arms (its variance spans 0.0--0.58 across seeds), so we do not report that metric. We include this as a negative result because a benchmark is well-served by an honest report that a fashionable architecture does not win on it (\texttt{results/gflownet\_seed\{0,1,2\}\_overnight.json}; a fair multi-seed diversity comparison at scale is future work, discussed in the future-work section).

\bigskip
\emph{Source: \texttt{results/\allowbreak prior\_\allowbreak vs\_\allowbreak learned.json}, \texttt{results/\allowbreak benchmark\_\allowbreak report\_\allowbreak depth2.json}, \texttt{results/\allowbreak benchmark\_\allowbreak report\_\allowbreak gap.json}, \texttt{docs/\allowbreak benchmark/\allowbreak stage2\_\allowbreak ranker\_\allowbreak evidence.md}, \texttt{results/\allowbreak selection\_\allowbreak ablation.json}, \texttt{results/\allowbreak selection\_\allowbreak ablation\_\allowbreak ranksignal.json}.}


\section{External validity of the coverage ceiling}
\label{app:external}
\subsection{Internal and external coverage ceilings}

The internal ceiling \texttt{coverage\_bank = 0.735} (the \textbf{micro}, pooled ratio-of-sums estimate $\Sigma_{\mathrm{hit}} / \Sigma_{\mathrm{true}}$ --- not a per-substrate mean; within-substrate pairs are dependent so a mean-of-ratios would misrepresent the uncertainty) carries a cluster-bootstrap 95\% CI of \textbf{[0.709, 0.762]} (resampling substrates, 10,000 resamples). Recomputed apples-to-apples on an \textbf{external} set --- the 37 GLORYx parent substrates, one uncapped full-bank depth-1 \texttt{apply\_rules} pass, no pool cap, the same tautomer match --- the external ceiling is \textbf{0.633} (95\% CI \textbf{[0.531, 0.733]}, $n=37$, wide by design at this sample size). This is far above the previously committed figure of \textbf{0.3715}, which is \textbf{pool-capped} --- a generator-budget artifact inherited from \texttt{gloryx\_oracle.json} --- and must \textbf{never} be read as ``the external ceiling.''

\subsection{Compositional analysis of the internal--external gap}

Internal (0.735) and external (0.633) both sit in the \textbf{micro} frame. A separate question is whether the internal-external gap is compositional: GLORYx parents tend to be larger, more-conjugated drugs than the internal test set, and the $\Delta$MW long tail noted earlier in the main text suggests coverage should degrade with molecular complexity. One OLS regression of \textbf{per-substrate} coverage on molecular descriptors (molecular weight, ring count, aromatic-atom count, heteroatom count, a conjugation-site count, and \texttt{n\_true}) fit across both populations predicts \textbf{macro} (per-substrate-mean) coverages of \textbf{$\sim$0.79} internal and \textbf{$\sim$0.74} external, in-sample --- the macro frame, distinct from the micro ceilings quoted above. Note honestly that this $\sim$0.74 in-sample external prediction sits \emph{above} the measured external macro coverage of 0.697, so the in-sample fit alone overstates transferability. Holding the external points fully out of the fit (fit-internal-only $\rightarrow$ predict-external) instead lands at \textbf{0.738 out-of-sample}, recovering \textbf{$\sim$56\%} of the internal$\rightarrow$external macro-coverage gap (\texttt{regression.predicted\_external\_mean\_oos = 0.738}, \texttt{regression.gap\_recovery\_frac = 0.565}). This out-of-sample number, not the in-sample one, is the transferable claim.

\subsection{Interpretation}

We frame this as a \textbf{suggestive partial composition effect at $n=37$} --- consistent with, but not proof of, molecular composition driving part of the internal-external coverage gap --- and explicitly \emph{not} a fitted or transferable law, and \emph{not} a defect in the rule bank or the matching protocol.

\emph{[Figure: internal-vs-external coverage scatter --- optional, post-draft]}

\emph{Source: \texttt{results/ceiling\_external\_validity.json}.}


\section{Released negative results}
\label{app:negative}
\subsection{0a. Rule-granularity probe (2026-07-19) --- measured, mostly negative}

Tested whether \textbf{rule granularity} (bank breadth/redundancy) is a lever on the recall factors \texttt{coverage(0.735) $\times$ selection(0.489) $\times$ ranking(0.726)}. Four measurements, each pre-registered with its falsifier; three of my own hypotheses were tempered or killed by the numbers.

\paragraph{Methodology caveat --- test peeked.}
These exploratory probes (\texttt{rule\_prune\_probe}, \texttt{ablate\_id\_embedding}, and the earlier budget-matched / cross-method runs) scored on the \textbf{test} split (\texttt{test\_predictions.csv}) --- a soft breach of \emph{touch-test-once}. Consequence: the final GRAIL-vs-MetaTox row is now slightly adaptively biased. Mitigation going forward: \textbf{test is frozen for that single final row} (spent on the GRAIL-vs-MetaTox comparison presented elsewhere in the paper); any follow-up exploratory run goes on \textbf{val}, and the val GRAIL predictions now exist (\texttt{artifacts/full5000\_single/predictions/val\_predictions.csv}, 994 substrates, produced by \texttt{dump\_val\_predictions.py} through the same code path as the test dump).

Verdict per claim, at measured confidence:

\begin{enumerate}
\item \textbf{Merging functional duplicates $\rightarrow$ SELECTION.} Redundancy is real: of the 2,747 rules that fire on a 1,000-substrate probe, 25\% are exact functional duplicates (\texttt{rule\_collapse.json}, held-out merge stability 0.975); a distribution-free canonical-SMIRKS dedup over all 7,581 removes $\sim$7.5\% (\texttt{rule\_dedup\_provable.json}; permutation variants are only $\sim$43\% of the redundancy, the rest is non-structural). \textbf{But the net selection benefit is unconfirmed:} \texttt{label\_reactions} marks every producing rule positive and each rule has its own \texttt{id\_embedding}, so duplicates split the training signal --- yet \texttt{noisy\_or} at inference pools the duplicate scores back (multiplicity $\approx$ frequency prior), partially healing the split. And the zero-id ablation (below) shows the split \texttt{id} is itself inert. \textbf{Verdict: weak lever; treat as deploy hygiene}, which folds into the resource+cache measurement (rule embeddings are substrate-independent and already cached once at startup --- bank size hits startup cost/memory, not per-query latency).

\item \textbf{Pruning dead rules.} \emph{The reducible-sparsity fact is solid; the selection mechanism moved, and the pruning scope shrank.} Training-wide (4,787 substrates, \texttt{rule\_train\_positives.json}): \textbf{56.3\% of the bank is NEVER a positive label} and \textbf{20.1\% is positive on exactly one substrate} $\rightarrow$ \textbf{$\sim$76\% of the TARGET is train-dead-or-singleton}, i.e. \emph{training} sparsity is reducible (P2-relevant, solid). Two cautions the self-audit forces, both previously over-stated:

\begin{itemize}
\item \textbf{Reachability (train/test confounder, third occurrence).} ``Prunable without coverage loss'' holds only for the \emph{training target} --- all-negative columns are pure train ballast. It is \textbf{not proven for the deployed bank}: a rule useless on train can be a true label on test/novel chemistry (e.g. a curated SyGMa general rule that never matched train). So ``76\% of the \emph{bank} prunes'' is unmeasured and carries \textbf{test-coverage risk}; only ``76\% of the \emph{target} is train-dead-or-singleton'' is solid.
\item \textbf{Mechanism (undercut by the same shot as merging).} The zero-id ablation shows \texttt{id} is inert, so ``dead rules inject noise via untrained \texttt{id}'' is \emph{also} inert --- the target-sparsity$\rightarrow$selection story is undermined. What survived was a \textbf{different, id-independent, inference-side} hypothesis: dead rules emit \textbf{distractor products} that dilute ranking, so cutting them could raise recall.
\item \textbf{RESOLVED --- and the lever is FALSIFIED (2026-07-19, \texttt{results/prune\_and\_rerank\_val.json}).} Pruning the 4,271 training-dead rules (bank 7,581 $\rightarrow$ 3,310) and re-ranking on \textbf{val} ($n=994$, tautomer, deployed operating point, prune set derived from TRAINING positives only) \textbf{hurts on both axes}: recall@15 \textbf{0.388 $\rightarrow$ 0.352} ($\Delta$ $-0.036$, 95\% CI $[-0.050, -0.024]$) and pool coverage \textbf{0.439 $\rightarrow$ 0.378} ($\Delta$ $-0.061$, CI $[-0.078, -0.046]$); mean pool 46.9 $\rightarrow$ 40.0. Both CIs exclude zero. So the pre-registered \textbf{reachability-loss} branch won, not distractor-cut: rules that never earn a positive label across 4,787 training substrates \textbf{still produce true metabolites on unseen val chemistry}. The reachability caveat above is now measured rather than hypothesised --- \emph{``76\% of the TARGET is train-dead-or-singleton'' stands; ``76\% of the BANK is prunable'' is FALSE.} The bank's apparent dead weight is generalization capacity living in its tail, which also argues against bank-shrinking as a deployment tactic (provable dedup remains safe; pruning does not). Note the ranking partially absorbs the damage --- 6.1 points of lost coverage cost only 3.6 points of realised recall --- but the net is unambiguously negative.
\end{itemize}

\item \textbf{\texttt{id\_embedding} dominance $\rightarrow$ mechanism of P2. FALSIFIED.} The rule embedding is \texttt{norm(graph\_encoding + id\_embedding + meta)}; the \texttt{id\_embedding} carries \textbf{82\% of cross-rule variance} vs 17\% for the GATv2 structure (\texttt{rule\_embed\_decomp.json}) --- seductively ``the model is a lookup table, rare-rule \texttt{id} is untrained noise, \emph{that} is P2.'' The \textbf{zero-id inference ablation kills it} (\texttt{ablate\_id\_embedding.json}): zeroing \texttt{id} moves recall@15 by $\mathbf{-0.007}$ and \emph{improves} recall@1/@5 ($+0.037$/$+0.027$). The 82\% variance is \textbf{not load-bearing} --- a geometric illusion; downstream scoring looks past \texttt{id}. P2 is not explained here.
\end{enumerate}

\paragraph{Byproduct --- VERIFIED ON VAL, and it does NOT replicate.}
The single test run had zeroing \texttt{id} \emph{raise} top-k precision ($+0.037$@1, $+0.027$@5, $n=291$), which I downgraded to ``anomaly, verify before believing''. Re-run on \textbf{val} with paired bootstrap CIs ($n=994$, 10k resamples, \texttt{results/ablate\_id\_embedding\_val.json}): \textbf{every $k$ is n.s.} --- @1 $+0.0011$ $[-0.0086, +0.0107]$, @5 $+0.0055$ $[-0.0074, +0.0181]$, @15 $-0.0116$ $[-0.0266, +0.0036]$, @30 $-0.0119$ $[-0.0282, +0.0044]$. The $+0.037$@1 was \textbf{single-run noise}; the ``free tweak'' is dead, exactly as the downgrade anticipated. What the val run \emph{does} confirm, now with CIs the original test run never had, is the main result: zeroing 82\% of the embedding variance moves recall by nothing distinguishable from zero at any $k$ --- \texttt{id} is \textbf{not load-bearing}, and the \texttt{id}$\rightarrow$P2 mechanism stays falsified on an independent split. \emph{(Method note: the first val attempt crashed at the CI step on a paired-vector alignment bug --- the per-substrate append was skipped for substrates with no candidates while \texttt{n} still counted them, so the two arms dropped different substrates (991 vs 990). It only crashed because the counts differed; equal counts would have produced a silently wrong CI. Fixed by recording recall 0 for empty candidates, plus an explicit alignment guard.)}

\paragraph{Net --- the track is CLOSED, and negatively.}
All three candidate levers were measured and none survived: \textbf{merging} = weak/deploy-hygiene (net effect unconfirmed; \texttt{noisy\_or} heals the training split), \textbf{pruning} = actively harmful (falsified on val: $-0.036$ recall, $-0.061$ coverage, both CIs excluding zero), \textbf{\texttt{id}$\rightarrow$P2} = falsified (zero-id ablation). What remains is one solid \emph{fact} (training sparsity is reducible --- 76\% of the \textbf{target} is train-dead-or-singleton, which is \emph{not} the same as the bank being prunable) and one structural insight (the bank's tail carries generalization capacity, so coverage really is the physical limit --- consistent with the main text's P1 diagnosis). Three of these outcomes falsified hypotheses I had argued for; that is the intended function of pre-registering the falsifier. Both deadline levers are now \textbf{done}: resource+cache (\texttt{resource\_cache\_profile.json}) and the GRAIL-vs-MetaTox comparison (presented elsewhere in the paper). Scripts: \texttt{prune\_and\_rerank\_val.py}, \texttt{rule\_collapse.py}, \texttt{rule\_dedup\_provable.py}, \texttt{rule\_prune\_probe.py}, \texttt{probe\_rule\_embeddings.py}, \texttt{ablate\_id\_embedding.py}.

\subsection{0c. Venue decision (2026-07-26) --- ICLR 2027 first, NeurIPS ED 2027 fallback}

\paragraph{Chosen:}
submit to \textbf{ICLR 2027} (abstract 19 Sep 2026, full paper \textbf{24 Sep 2026} --- $\sim$8 weeks), with \textbf{NeurIPS Evaluations \& Datasets 2027} ($\sim$May 2027) as the fallback. Sequencing works: an ICLR decision lands $\sim$Jan 2027, leaving time to fold the reviews into an ED submission.

\emph{Why not the ideal venue first:} NeurIPS 2026's ED track (the renamed D\&B track) closed \textbf{6 May 2026}, $\sim$2.5 months before the manuscript was ready. Waiting for ED 2027 costs $\sim$10 months in a neighbourhood the lit-check found to be actively standardizing (\citet{liu2026massspecgymwild}, retrosynthesis-benchmark rethinking 2026, \citet{agarwal2026metric}) --- real scoop risk.

\emph{Risk that made this a close call, and what defuses it:} an ML main track judges by method norms, where ``our own instrument loses on recall'' reads as a defect rather than a design principle. \textbf{ICLR carries a \texttt{Datasets and benchmarks} primary area}, so the submission routes to reviewers who expect a benchmark contribution --- this largely restores the category protection the external review said D\&B venues provide. Journals (JCIM / J. Cheminformatics, Q1, rolling) were the low-risk alternative and were declined in favour of the A* attempt; no institutional constraint forced a journal.

\paragraph{Conversion work (the real cost).}
ICLR is \textbf{9 pages main text + unlimited appendix, LaTeX}; the current draft is $\sim$13.5k words of markdown, so roughly two thirds moves to the appendix. Proposed split --- \emph{main:} the two failure modes, TAME, the rank-flip (2 internal pairs \textbf{+ the external GLORYx-37 replication}, now the strongest single piece of evidence), the coverage $\times$ selection $\times$ ranking decomposition, P2-as-XMC in brief, the honest anchor, limitations. \emph{Appendix:} full architecture (main-text architecture section), formal framework (main-text formal-framework section), proposition witnesses, the rule-granularity negative results, the GFlowNet negative result, external-validity regression, cross-domain probe.

\paragraph{Practical constraints to handle:}
(i) ICLR is \textbf{double-blind} and the repository is public (\texttt{github.com/doctawho42/GRAIL}) --- an anonymized artifact link is required at submission; (ii) the 275-word abstract needs no trimming (ICLR has no JCIM-style 250-word cap); (iii) the MetaTox decision (covered elsewhere in the paper) resolves by deadline --- include as a 6th method only if the SMIRKS variant arrives and is re-scored in time, otherwise keep the leaderboard at five and cite the comparison as external validation.


\section{Limitations and scope}
\label{app:xdomain}
\subsection{No recall win, and precision as a pessimistic lower bound}

We state these plainly rather than let them surface in review. First, \textbf{no recall win}: GRAIL does not win on recall anywhere in this paper. It loses to SyGMa by a certified, significant margin (the certification section: paired $\Delta = -0.242$, 95\% CI $[-0.271, -0.212]$; McNemar $p \approx 1.7\times10^{-44}$) and to the learned transformer baselines on the shared $n=150$ subset (MetaPredictor 0.585, MetaTrans 0.561, both above GRAIL's 0.365 under the same tautomer protocol, the match-sensitivity section). Nothing in this paper should be read as a state-of-the-art claim; the contribution is the ceiling, the decomposition, and the protocol, with GRAIL as one honestly diagnosed row.

Second, \textbf{precision is a pessimistic lower bound} throughout: because negatives are rule-applicable, non-annotated products rather than confirmed non-metabolites (the data-labeling section), an unannotated true metabolite is scored as a false positive, so every precision number in this paper understates the pipeline's true precision by an unknown, non-annotation-corrected amount --- we lead with recall and \texttt{mean\_output\_size} for this reason and do not report precision as a headline metric.

\subsection{Single-checkpoint headline and a mixed-sample-size comparison}

Third, the deployed headline (recall@15 $= 0.261$ micro / $0.330$ macro, the headline-results and certification sections) is a \textbf{single checkpoint}, not a seed-averaged estimate; the honest-anchor certification (the certification section) bounds \emph{evaluation} variance --- resampling and discordance over a fixed model --- not \emph{training} variance across independently seeded runs, and a three-seed retraining confirms the headline is seed-stable (macro $0.344 \pm 0.010$, micro $0.269 \pm 0.006$; the certification section).

Fourth, the five-method, five-protocol match-sensitivity comparison (the match-sensitivity section, Table 3) mixes sample sizes: the three tier-2 comparators (BioTransformer, MetaPredictor, MetaTrans) are scored only on the \textbf{$n=150$} shared subset for which frozen third-party predictions exist, while GRAIL and SyGMa's headline numbers elsewhere in the paper are reported on the full \textbf{$n\approx1170$} clean test split; within Table 3 itself GRAIL and SyGMa are re-scored on the matching $n=150$ subset so the five columns are paired on identical substrates. MetaPredictor has since been re-scored on the full test (recall@15 $0.568$, $\approx$ SyGMa, \texttt{results/metapredictor\_1170.json}), confirming its $n=150$ value is representative and leaving only BioTransformer and MetaTrans on $n=150$; a full single-$n$, five-method rerun remains future work (BioTransformer's exact $n=150$ configuration is unscripted and MetaTrans's inference pipeline is no longer reproducible in-tree).

\subsection{External validity, an unmeasured identifiability assumption, and single-step scope}

Fifth, the external-validity ceiling (the external-validity section) is measured on only \textbf{37} GLORYx parent substrates, so its 95\% CI is wide by design ($[0.531, 0.733]$) and the composition-effect regression that partially explains the internal--external gap is fit on the same small external population --- we frame it as suggestive, not as a transferable law.

Sixth, Proposition 2's explanatory model rests on an \textbf{unmeasured assumption}: the propensity model $e(r) \propto \pi(r)$ (labeling probability proportional to the marginal rule-firing rate) is asserted, not estimated from data, so the identifiability argument for why the learned selector underperforms a frequency prior is consistent with, but not proof of, the stated mechanism (the propositions-and-mechanism section).

Finally, Proposition 3's paradigm bound is \textbf{single-step conditional}: depth-2 rule application recovers only $+0.012$ ceiling at $8.5\times$ candidate cost (the propositions-and-mechanism section), so multi-step and out-of-bank chemistry remain an open, unresolved coverage lever rather than a closed question --- the bound constrains the single-step paradigm, not the metabolite-prediction problem in general.

\subsection{Factorized-generator intervention}

We also built and evaluated a factorized-generator redesign of Stage-1 as a direct intervention (the propositions-and-mechanism section); it does \textbf{not} beat the deployed pipeline as a generator (coverage-bound at recall@15 $0.256 <$ deployed $0.365$) --- its only deployable value is a modest, paired-significant re-ranking gain ($+0.0165$, 95\% CI $[0.006, 0.027]$, full clean test, $n=1170$) --- and closing the larger gap to SyGMa still requires a broader rule bank, out of scope here, consistent with the coverage-binding finding.

\subsection{Scope and generality: a cross-domain protocol probe}

Finally, a \textbf{scope-and-generality} limitation that doubles as this paper's clearest future direction: our rank-flip result (the match-sensitivity section) is established \textbf{within metabolite prediction only}. Whether match-protocol choice reorders leaderboards \emph{across} molecular-ML tasks --- retrosynthesis, molecular generation --- is an open, testable hypothesis, and we name its falsifying run explicitly: re-score $\geq 2$ comparable published models' \emph{frozen} predictions in another domain under the same five match conventions; if the method ordering does not move (interaction CI covering zero), the generality claim fails.

A clean instance requires $\geq 2$ models of \textbf{comparable strength but differing match-sensitivity} whose raw ranked predictions are public; we found they are not (only same-family or strength-mismatched checkpoints are readily runnable), so the multi-model cross-domain rank-flip is deferred rather than forced.

A first single-model probe (ReactionT5v2 on the USPTO-50k test, $n=200$; \texttt{results/xdomain\_retro\_protocol.json}) in fact suggests the effect is \textbf{domain-dependent, not universal}: the \emph{same} predictions move only \textbf{$\sim 0.02$} in top-1 accuracy across \{canonical, stereo-stripped, InChIKey, tautomer\} matching ($0.695$--$0.715$) --- roughly an order of magnitude below the metabolism interaction (the match-sensitivity section) --- plausibly because retrosynthesis references are canonical-clean small molecules, whereas the metabolism rule engine emits tautomer/stereo variants of its references.

So the honest cross-domain reading is \textbf{not} that protocol choice universally reorders leaderboards, but that its magnitude is a function of how much a generator's outputs diverge tautomerically/stereochemically from the references. Establishing whether protocol choice \emph{reorders} (not merely shifts) leaderboards in a domain with comparable tautomer ambiguity remains the named open test, and needs only public multi-model predictions in a re-scorable form --- which do not yet exist.


\section{Worked example: one prediction set, two verdicts}
\label{app:worked}
\textbf{Worked example.} The quotient dependence is not hypothetical. Consider a predicted set \texttt{\{D-alanine, acetone (enol tautomer)\}} scored against annotated references \texttt{\{L-alanine, acetone (keto tautomer)\}}. Under strict \texttt{inchikey}, both predictions miss: the alanine pair differs at a stereocenter, and the acetone pair differs only by a keto/enol tautomerization outside the subset plain InChI normalizes, so it also misses under \texttt{inchi\_no\_stereo}. \texttt{inchikey\_tautomer} standardization additionally strips stereochemistry (canonical SMILES generated with \texttt{isomericSmiles=False}), so tautomer-canonicalizing both sides collapses both gaps at once: the acetone keto/enol pair becomes a hit \emph{and} the D-/L-alanine stereocenter pair becomes a hit, scoring full recall on this example. The same fixed predictions can therefore score as a near-total miss or a full hit purely by changing the match quotient, with no change to the underlying chemistry --- the phenomenon quantified across methods and the shared external set in the matching-quotient sensitivity analysis.


\section{Budget-matched comparison}
\label{app:budget}
Recall at a fixed cut-off is not a comparison at matched cost. In their deployed configurations
SyGMa reaches 0.554 by emitting about 74 predictions per substrate; GRAIL reaches 0.365 at about
8.7. Normalising by output size reorders the table: SyGMa's precision is 0.029, roughly four times
below the low-output methods: GRAIL at 0.109, BioTransformer at 0.127, MetaTrans at 0.122 and
MetaPredictor at 0.116. Precision understates all of these under incomplete annotation, but the
sevenfold difference in output volume it reflects is real. At a smaller matched budget the recall
spread compresses as well: at $k=5$, SyGMa's 0.465 is close to MetaPredictor's 0.470, with GRAIL at
0.321. This matches the independent finding of \citet{Boyce_2022} that SyGMa buys coverage with
prediction volume.

\begin{table}[h]
\centering
\caption{Recall at matched cut-offs, precision, and mean output size under tautomer-aware
matching, for each method in its deployed configuration on the shared subset ($n=\nshared$). Recall
here is \emph{macro}, the mean of per-substrate recalls; the decomposition below reports the same
methods on the same substrates in \emph{micro}, the pooled ratio, which is why GRAIL reads $0.365$
in this table and $0.318$ there.}
\label{tab:budget-matched-leaderboard}
\begin{tabular}{lccccc}
\toprule
Method & Recall@5 & Recall@10 & Recall@15 & Precision & Mean output \\
\midrule
GRAIL & 0.321 & 0.356 & 0.365 & 0.109 & 8.7 \\
SyGMa & 0.465 & 0.533 & 0.554 & 0.029 & 74.2 \\
BioTransformer & 0.315 & 0.396 & 0.444 & 0.127 & 10.8 \\
MetaPredictor & 0.470 & 0.575 & 0.585 & 0.116 & 11.2 \\
MetaTrans & 0.424 & 0.533 & 0.561 & 0.122 & 12.7 \\
\bottomrule
\end{tabular}
\end{table}

\subsection{The frontier across budgets}

Table~\ref{tab:budget-matched-leaderboard} compares methods as they are deployed, each with its own
output policy. A second measurement removes that difference: rank each method's own candidate pool
by its own score and truncate at a common $k$. This is a separate re-scoring run on the same
$\nshared$ substrates and the same matching criterion, and its numbers are not the deployed ones,
because truncating a pool at $k$ is not what either system does in normal use.

Under that common budget, measured at seven cut-offs between 1 and 64, SyGMa leads GRAIL at every
one with no crossover: 0.233 against 0.127 at $k=1$, 0.493 against 0.350 at $k=15$, and 0.520
against 0.410 at $k=64$. GRAIL does not win at a matched output budget, and its selector and ranker
are the binding weakness.

\subsection{Decomposing both methods}

Applying the decomposition to each method, rather than to GRAIL alone, shows the mechanism. The
values below are not the frontier's, and they are \emph{micro} where the table above is macro, so
GRAIL reads $0.318$ here against $0.365$ there on the same substrates. They describe each method's
\emph{deployed} output on the same
$\nshared$ substrates under the same criterion, whereas the frontier truncates each method's full
candidate pool at a common $k$. The two differ for GRAIL, which reports $0.318$ as deployed against
$0.350$ at $k=15$ on the frontier, because the frontier draws fifteen candidates from a pool the
deployed selector never emits.

SyGMa has no selection stage: it applies its whole applicable rule set and emits the result, at
about 64 outputs per substrate on this subset. Almost nothing is lost between what its pool
reaches and what it reports, 0.520 against a recall@15 of 0.493. Its recall is therefore close to
the coverage its bank realises.

GRAIL's bank reaches more of these transformations than SyGMa emits: its ceiling on this subset is
0.701, against a full-test ceiling of $\ceiling$ and SyGMa's realised 0.520. GRAIL nevertheless
reports 0.318, because selection and ranking together retain 0.454 of what the bank could reach.
Nothing is lost at the cut-off itself \emph{on these $\nshared$ substrates}: emitting about eight
candidates on average, GRAIL's budgeted pool is not truncated here, so its recall@15 equals that
pool's coverage exactly. That is a property of this subset and not of the method: on the full split
GRAIL's output policy, which is itself a cap at fifteen, is reached on $308$ of $\ntest$ substrates, and
that is where the third factor of \S\ref{sec:decomp} loses its references.

The gap to SyGMa is thus not a worse rule bank. It is the cost of having a selection stage at all,
which recall charges only to methods that select, and which GRAIL currently pays without
recovering.


\end{document}
